\section{Trace four frames through the averaging coefficients}
Set all eight coefficient codes to 4096. Each represents one eighth, so the exact
sum followed by the single prescribed rounding computes an eight-position moving
average with zero beginning history. The first four frames never fill the window;
its missing older values remain zeros, rather than reducing the divisor.

\tagpdfsetup{table/header-rows={1}}
\begin{center}
  \begin{tabular}{rrrrr}
    Frame & Input left & Input right & Output left & Output right \\
    \hline
    0     & 12000      & $-12000$    & 1500        & $-1500$      \\
    1     & 0          & 16000       & 1500        & 500          \\
    2     & $-8000$    & 0           & 500         & 500          \\
    3     & 32767      & $-32768$    & 4596        & $-3596$      \\
  \end{tabular}
\end{center}

At frame three, the left window's four nonzero-time inputs sum to 36767. Dividing
by eight gives 4595.875, which rounds to 4596. The right sum is minus 28768, giving
minus 3596 exactly. No result clips. In interleaved little-endian storage the sixteen
result bytes are \texttt{dc 05 24 fa dc 05 f4 01 f4 01 f4 01 f4 11 f4 f1}.

These results establish the beginning condition, channel separation, rounding and
stored byte order together. They do not establish clipping behavior: other coefficient
sets can amplify a window far beyond either endpoint. Keep saturation and tie cases
as separate tests, even though an averaging example sounds and looks well behaved.
